\scanpage{117}
% archive foliation: 116
D'autre part, faisant $\unsure{t=r|t=s=r}$ dans (2), on trouve
\[
G(t,t)=\deleted{\unsure{G}}0\qquad(t\text{ uniformisant})
\]
et faisant $r=t$, on trouve
\[
G(t,s)+G(s,t)=0,\qquad(s,t\text{ uniformisants})
\]
\inserted{i.e. $G(t,s)=-G(s,t)$}
\deleted{\unsure{donc on} prolonge l'expression $G(t,s)$, définie dans le voisinage par ($t:X\longrightarrow\mathbb E^1(\mathbb C)$, $s\longrightarrow\mathbb P^1(\mathbb C)$), au cas où $t$ peut aussi admettre des pôles.}
\margin{\unsure{ceci permet} \unsure{par (1)} \illegible{2} \unsure{facilement pour} $s,t:X\longrightarrow\mathbb P^1(\mathbb C)$.}
d'où
\[
(1\mathrm{bis})\quad\left\{\begin{array}{l}
G(at,s)=G(t,s)\\
G(t+b,s)=G(t,s)\\
G(\frac1t,s)=G(t,s)\qquad(t,\frac1t\text{ unif}^{\mathrm{ts}})
\end{array}\right.
\]
ce qui permet de faire la $=$ généralisation pour les $t$ que pour les $s$, d'où pour les
\[
s,t:X\longrightarrow\mathbb P^1(\mathbb C)\qquad\text{morphismes étales}
\]
\[
(4)\quad\left\{\begin{array}{l}
G(\frac{at+b}{ct+d},s)=G(t,s)\qquad a,b,c,d\in\mathbb C\text{ avec }ad-bc\in\mathbb C^*\\
\text{en plus de}\\
G(t,\frac{as+b}{cs+d})=G(t,s)\qquad\unsure{''}
\end{array}\right.
\]

\deleted{De plus on a}
Soit $\mathcal P_X$ le faisceau des morphismes \inserted{analytiques} étales
\[
X\longrightarrow\mathbb P^1(\mathbb C),
\]
sur lequel opère le groupe
\[
H=\operatorname{Aut}_{\mathbb C}\mathbb P^1(\mathbb C)\simeq\mathrm{GP}(1)(\mathbb C):=\left\{\left(\begin{matrix}a&b\\c&d\end{matrix}\right)\,\middle|\,ad-bc\in\mathbb C^*\right\}\quad\text{mod. matrices diagonales}.
\]
On a donc défini un morphisme de faisceaux
\[
\underline G:\mathcal P_X\times\mathcal P_X\longrightarrow\deleted{\unsure{\Omega}}\underline\omega_X^{\otimes2}
\]
satisfaisant les deux conditions

I) $G$ invariant par $H\times H$ i.e.
\[
G(gt,g's)=G(t,s)\qquad\text{si }g,g'\in H
\]
II) $G(t,s)+G(s,r)=G(t,r)$

En particulier $G(t,t)=0$, $G(t,s)=-G(s,t)$ \unsure{(pour $s$, pour $t$ fixé)}.

En fait, la théorie des équations différentielles nous donne aisément une réciproque à I), savoir
